%! TeX program = pdflatex
\documentclass[aspectratio=169,12pt]{beamer}

\usepackage[T1]{fontenc}
\usepackage[default]{sourcesanspro}
\usefonttheme[onlymath]{serif}
\usepackage{amsmath,amssymb,mathtools}
\usepackage{booktabs}
\usepackage{array}
\usepackage{microtype}
\usepackage{tikz}
\usetikzlibrary{
  arrows.meta,
  backgrounds,
  calc,
  decorations.pathreplacing,
  fit,
  positioning,
  shapes.geometric
}
\usepackage{beamerthemestochastic}

\newcommand{\OPT}{\mathrm{OPT}}
\newcommand{\optpost}{\OPT_{\mathrm{post}}}
\newcommand{\optadapt}{\OPT_{\mathrm{adapt}}}
\newcommand{\optprior}{\OPT_{\mathrm{prior}}}
\newcommand{\E}{\mathbb{E}}
\newcommand{\Pp}{\mathbb{P}}
\newcommand{\cost}{\operatorname{cost}}
\newcommand{\poly}{\operatorname{poly}}
\newcommand{\1}{\mathbf{1}}

\title[Information Gaps \& Algorithms]{%
  Information Gaps and Algorithms for the\\
  Stochastic Traveling Salesperson Problem}
\author{Teng Liu}
\institute{Department of Computer Science \enspace·\enspace ETH Zürich}
\date{Master Thesis \enspace·\enspace July 2026}

% Native TikZ figures for the thesis talk.
% All figures inherit the palette from beamerthemestochastic.sty.

\tikzset{
  stDet/.style={
    rectangle,
    rounded corners=0.45pt,
    draw=STInk,
    fill=STWhite,
    line width=0.7pt,
    minimum width=0.52cm,
    minimum height=0.48cm,
    inner sep=1.5pt,
    font=\fontsize{8.5pt}{9pt}\selectfont
  },
  stDepot/.style={
    stDet,
    fill=STGold!25,
    draw=STGold!65!STInk,
    line width=0.9pt
  },
  stStoch/.style={
    circle,
    draw=STInk,
    fill=STTeal!20,
    line width=0.7pt,
    minimum size=0.47cm,
    inner sep=1pt,
    font=\fontsize{8pt}{9pt}\selectfont
  },
  stStochActive/.style={
    stStoch,
    fill=STTeal,
    text=STWhite,
    draw=STTeal!55!STInk
  },
  stStochInactive/.style={
    stStoch,
    fill=STPaper,
    draw=STMuted,
    text=STMuted
  },
  stArrow/.style={
    -{Stealth[length=2.2mm,width=1.45mm]},
    line width=1.05pt,
    draw=STBlue
  },
  stArrowThin/.style={
    -{Stealth[length=1.8mm,width=1.2mm]},
    line width=0.72pt,
    draw=STBlue
  },
  stEdge/.style={draw=STInkSoft,line width=0.70pt},
  stFaintEdge/.style={draw=STLine,line width=0.5pt},
  stDashed/.style={draw=STMuted,densely dashed,line width=0.7pt},
  stLabel/.style={
    font=\fontsize{8.5pt}{10pt}\selectfont,
    text=STMuted,
    align=center
  },
  stSmallLabel/.style={
    font=\fontsize{7.5pt}{9pt}\selectfont,
    text=STMuted,
    align=center
  }
}

% --------------------------------------------------------------------------
% Background: a master tour and a realized shortcut
% --------------------------------------------------------------------------
\newcommand{\StMasterOrder}{%
\begin{tikzpicture}[x=1cm,y=1cm]
  \node[anchor=west,font=\fontsize{9pt}{11pt}\selectfont\bfseries,text=STInk]
    at (0,2.65) {Master order};
  \node[stDepot] (mr) at (0.7,2.0) {$r$};
  \node[stStochActive] (mone) at (2.0,2.0) {$1$};
  \node[stStochInactive] (mtwo) at (3.3,2.0) {$2$};
  \node[stStochActive] (mthree) at (4.6,2.0) {$3$};
  \node[stStochInactive] (mfour) at (5.9,2.0) {$4$};
  \node[stStochInactive] (mfive) at (7.2,2.0) {$5$};
  \node[stStochActive] (msix) at (8.5,2.0) {$6$};
  \node[stDepot] (mrb) at (9.8,2.0) {$r$};
  \foreach \a/\b in {r/one,one/two,two/three,three/four,four/five,five/six,six/rb}{
    \draw[stFaintEdge,-{Stealth[length=1.4mm]}] (m\a)--(m\b);
  }
  \node[stSmallLabel,anchor=north] at (3.3,1.63) {inactive};
  \node[stSmallLabel,anchor=north] at (5.9,1.63) {inactive};
  \node[stSmallLabel,anchor=north] at (7.2,1.63) {inactive};

  \node[anchor=west,font=\fontsize{9pt}{11pt}\selectfont\bfseries,text=STInk]
    at (0,0.93) {Realized tour};
  \node[stDepot] (rr) at (0.7,0.28) {$r$};
  \node[stStochActive] (a1) at (2.8,0.28) {$1$};
  \node[stStochActive] (a3) at (5.15,0.28) {$3$};
  \node[stStochActive] (a6) at (7.55,0.28) {$6$};
  \node[stDepot] (rr2) at (9.8,0.28) {$r$};
  \draw[stArrow] (rr)--node[above,stSmallLabel] {$d(r,1)$} (a1);
  \draw[stArrow] (a1)--node[above,stSmallLabel] {$d(1,3)$} (a3);
  \draw[stArrow] (a3)--node[above,stSmallLabel] {$d(3,6)$} (a6);
  \draw[stArrow] (a6)--(rr2);
  \node[stLabel,anchor=west,text=STTeal] at (10.25,0.28)
    {inactive clients\\are shortcut};
\end{tikzpicture}}

% --------------------------------------------------------------------------
% Background: three information regimes
% --------------------------------------------------------------------------
\newcommand{\StBenchmarkLadder}{%
\begin{tikzpicture}[x=1cm,y=1cm]
  \draw[STLine,line width=2.0pt] (0.7,0)--(11.1,0);
  \draw[STTeal,line width=2.8pt] (0.7,0)--(3.55,0);
  \draw[STBlue,line width=2.8pt] (3.55,0)--(7.25,0);
  \draw[STGold,line width=2.8pt] (7.25,0)--(11.1,0);

  \foreach \x/\col/\num in {1.05/STTeal/1,5.4/STBlue/2,10.5/STGold/3}{
    \node[circle,fill=\col,text=STWhite,minimum size=0.52cm,
      font=\fontsize{10pt}{11pt}\selectfont\bfseries] at (\x,0) {\num};
  }

  \node[anchor=north west,text width=3.05cm,align=left] at (0.55,2.70) {%
    {\fontsize{14pt}{16pt}\selectfont\bfseries A posteriori}\par
    \vspace{0.12cm}
    {\fontsize{10.5pt}{13pt}\selectfont\color{STInkSoft}
      Know the active set\\before choosing.}};
  \node[anchor=north,text width=3.55cm,align=center] at (5.4,2.70) {%
    {\fontsize{14pt}{16pt}\selectfont\bfseries Adaptive}\par
    \vspace{0.12cm}
    {\fontsize{10.5pt}{13pt}\selectfont\color{STInkSoft}
      Learn on calls;\\move when active.}};
  \node[anchor=north east,text width=3.35cm,align=right] at (11.0,2.70) {%
    {\fontsize{14pt}{16pt}\selectfont\bfseries A priori}\par
    \vspace{0.12cm}
    {\fontsize{10.5pt}{13pt}\selectfont\color{STInkSoft}
      Fix one\\master order.}};

  \node[anchor=north west,text=STTeal,
    font=\fontsize{9.5pt}{11.5pt}\selectfont\bfseries] at (0.62,-0.38)
    {knows the active set};
  \node[anchor=north,text=STBlue,
    font=\fontsize{9.5pt}{11.5pt}\selectfont\bfseries] at (5.4,-0.38)
    {learns causally};
  \node[anchor=north east,text=STGold!75!STInk,
    font=\fontsize{9.5pt}{11.5pt}\selectfont\bfseries] at (11.0,-0.38)
    {does not observe the set};
\end{tikzpicture}}

% --------------------------------------------------------------------------
% Construction I: cycle-subdivided clique
% --------------------------------------------------------------------------
\newcommand{\StCycleClique}{%
\begin{tikzpicture}[scale=0.90]
  \foreach \i in {1,...,6} {
    \coordinate (c\i) at ({90-60*(\i-1)}:3.0);
  }
  \foreach \i in {1,...,6} {
    \pgfmathtruncatemacro{\next}{mod(\i,6)+1}
    \coordinate (m\i) at ($(c\i)!0.5!(c\next)$);
  }
  \foreach \i in {1,...,6} {
    \foreach \j in {1,...,6} {
      \ifnum\j>\i
        \pgfmathtruncatemacro{\next}{mod(\i,6)+1}
        \ifnum\j=\next\else
          \ifnum\i=1
            \ifnum\j=6\else
              \draw[STLine,line width=0.5pt] (c\i)--(c\j);
            \fi
          \else
            \draw[STLine,line width=0.5pt] (c\i)--(c\j);
          \fi
        \fi
      \fi
    }
  }
  \foreach \i in {1,...,6} {
    \pgfmathtruncatemacro{\next}{mod(\i,6)+1}
    \draw[STBlue,line width=1.55pt] (c\i)--(m\i)--(c\next);
  }
  \foreach \i in {1,...,6} {
    \node[stDet,minimum width=0.62cm,minimum height=0.57cm] (v\i) at (c\i)
      {$v_{\i}$};
    \node[stStoch,minimum size=0.53cm] at (m\i) {$v'_{\i}$};
  }
  \node[stSmallLabel,anchor=south] at ([yshift=0.36cm]v1.north)
    {depot};
  \node[stSmallLabel,text=STBlue,anchor=north] at (0,-3.46)
    {thick cycle edges are subdivided};
\end{tikzpicture}}

\newcommand{\StSegmentTax}{%
\begin{tikzpicture}[x=1cm,y=1cm]
  \node[stDet,minimum size=0.60cm] (u) at (0.35,0) {$u$};
  \node[stStoch] (w1) at (1.65,0) {$w_1$};
  \node[stStoch] (w2) at (2.75,0) {$w_2$};
  \node[font=\fontsize{14pt}{15pt}\selectfont,text=STMuted] at (3.75,0) {$\cdots$};
  \node[stStoch] (wk) at (4.75,0) {$w_k$};
  \node[stDet,minimum size=0.60cm] (v) at (6.1,0) {$v$};
  \draw[stEdge] (u)--(w1)--(w2);
  \draw[stEdge] (w2)--(3.43,0);
  \draw[stEdge] (4.07,0)--(wk)--(v);
  \draw[decorate,decoration={brace,amplitude=5pt,mirror},STBlue,line width=0.8pt]
    (0.05,-0.58)--(6.4,-0.58)
    node[midway,below=0.25cm,text=STBlue,
      font=\fontsize{10.5pt}{12pt}\selectfont\bfseries]
      {$\Phi(u;W;v)\ge 1+\frac{7}{8}k$};

  \node[anchor=south,text=STInkSoft,
    font=\fontsize{9.2pt}{11pt}\selectfont] at (3.22,0.62)
    {one segment with \(k\) independent subdividers};
\end{tikzpicture}}

\newcommand{\StAdaptivePeel}{%
\begin{tikzpicture}[x=1cm,y=1cm]
  \node[stDet] (l) at (0.3,0) {};
  \node[stStochActive] (a) at (1.2,0) {};
  \node[stDet] (u) at (2.1,0) {};
  \node[stStoch] (b) at (3.0,0) {};
  \node[stDet] (c) at (3.9,0) {};
  \node[stStochInactive] (d) at (4.8,0) {};
  \node[stDet] (r) at (5.7,0) {};
  \draw[STLine,line width=1.1pt] (l)--(a)--(u)--(b)--(c)--(d)--(r);
  \draw[stArrow,draw=STTeal] ([yshift=0.15cm]l.north east)
    to[bend left=18] node[above=0.04cm,stSmallLabel,text=STTeal]
    {active: traverse} ([yshift=0.13cm]u.north west);
  \draw[stArrow,draw=STCoral] ([yshift=-0.12cm]c.south east)
    to[bend left=43] node[below=0.12cm,stSmallLabel,text=STCoral]
    {inactive: switch endpoint} ([yshift=-0.12cm]l.south west);
  \node[stSmallLabel,anchor=north] at (4.8,-0.38) {skip};
  \draw[decorate,decoration={brace,amplitude=4pt},STBlue,line width=0.7pt]
    (0.05,0.88)--(5.95,0.88)
    node[midway,above=0.18cm,stLabel,text=STBlue]
    {uncalled clients remain one interval};
\end{tikzpicture}}

% --------------------------------------------------------------------------
% Construction II: three chains and incompatible skeletons
% --------------------------------------------------------------------------
\newcommand{\StThreeChain}{%
\begin{tikzpicture}[x=0.95cm,y=0.95cm]
  \node[stDepot] (A0) at (0, 2.35) {$a_0$};
  \node[stDet] (A1) at (1.4, 2.35) {$a_1$};
  \node[stDet] (Ah) at (6.0, 2.35) {$a_h$};
  \node[stDet] (Bh) at (0, 0) {$b_h$};
  \node[stDet] (B1) at (4.6, 0) {$b_1$};
  \node[stDet] (B0) at (6.0, 0) {$b_0$};
  \node[stDet] (C0) at (0,-2.35) {$c_0$};
  \node[stDet] (C1) at (1.4,-2.35) {$c_1$};
  \node[stDet] (Ch) at (6.0,-2.35) {$c_h$};
  \node[stStoch,minimum size=0.60cm] (S) at (7.35,-1.18) {$s$};

  \draw[stDashed] (A0)--(Bh);
  \draw[stDashed] (A0)--(Ch);
  \draw[stDashed] (B0)--(Ah);
  \draw[stDashed] (C0)--(Ah);
  \draw[stDashed] (C0)--(Bh);
  \draw[stDashed] (B0)--(S);
  \draw[stDashed] (Ch)--(S);

  \draw[STBlue,line width=1.15pt] (A0)--(A1)--(Ah);
  \draw[STBlue,line width=1.15pt] (Bh)--(B1)--(B0);
  \draw[STBlue,line width=1.15pt] (C0)--(C1)--(Ch);
  \node[fill=STPaper,inner sep=3pt,text=STMuted] at (3.6,2.35) {$\cdots$};
  \node[fill=STPaper,inner sep=3pt,text=STMuted] at (2.5,0) {$\cdots$};
  \node[fill=STPaper,inner sep=3pt,text=STMuted] at (3.6,-2.35) {$\cdots$};
  \node[stLabel,anchor=east,font=\fontsize{10pt}{12pt}\selectfont\bfseries]
    at (-0.55,2.35) {$A$};
  \node[stLabel,anchor=east,font=\fontsize{10pt}{12pt}\selectfont\bfseries]
    at (-0.55,0) {$B$};
  \node[stLabel,anchor=east,font=\fontsize{10pt}{12pt}\selectfont\bfseries]
    at (-0.55,-2.35) {$C$};
  \node[stSmallLabel,anchor=west,text=STBlue] at (7.0,1.45)
    {chain edges: $1$};
  \node[stSmallLabel,anchor=west,text=STMuted] at (7.0,0.95)
    {dashed: $2h$};
\end{tikzpicture}}

\newcommand{\StSkeletonComparison}{%
\begin{tikzpicture}[x=1cm,y=1cm]
  % Left panel: s inactive.
  \node[anchor=west,font=\fontsize{12pt}{14pt}\selectfont\bfseries,text=STBlue]
    at (0,2.25) {$s$ inactive};
  \node[stDet] (A0) at (0.45,0.85) {$A$};
  \node[stDet] (B0) at (2.05,0.85) {$B$};
  \node[stDet] (C0) at (3.65,0.85) {$C$};
  \draw[stArrow,draw=STBlue] (A0)--(B0);
  \draw[stArrow,draw=STBlue] (B0)--(C0);
  \draw[stArrow,draw=STBlue] (C0) to[bend left=38] (A0);
  \node[stLabel,text=STBlue] at (2.05,-0.08) {three crossings};
  \node[font=\fontsize{17pt}{19pt}\selectfont\bfseries,text=STInk]
    at (2.05,-0.83) {$9h$};

  \draw[STLine,line width=0.7pt] (5.15,-1.2)--(5.15,2.35);

  % Right panel: s active.
  \node[anchor=west,font=\fontsize{12pt}{14pt}\selectfont\bfseries,text=STCoral]
    at (5.65,2.25) {$s$ active};
  \node[stDet] (A1) at (6.1,0.85) {$A$};
  \node[stDet] (C1) at (7.55,0.85) {$C$};
  \node[stStochActive] (S1) at (9.0,0.85) {$s$};
  \node[stDet] (B1) at (10.45,0.85) {$B$};
  \draw[stArrow,draw=STCoral] (A1)--(C1);
  \draw[stArrow,draw=STCoral] (C1)--(S1);
  \draw[stArrow,draw=STCoral] (S1)--(B1);
  \draw[stArrow,draw=STCoral] (B1) to[bend left=35] (A1);
  \node[stLabel,text=STCoral] at (8.28,-0.08) {four crossings};
  \node[font=\fontsize{17pt}{19pt}\selectfont\bfseries,text=STInk]
    at (8.28,-0.83) {$11h$};
\end{tikzpicture}}

% --------------------------------------------------------------------------
% Construction III: recursive directed block
% --------------------------------------------------------------------------
\newcommand{\StRecursiveBlock}{%
\begin{tikzpicture}[x=1cm,y=1cm]
  \node[draw=STBlue,dashed,rounded corners=0.65pt,minimum width=2.35cm,
    minimum height=2.15cm,fill=STBlue!4] (L) at (1.35,0) {};
  \node[draw=STBlue,dashed,rounded corners=0.65pt,minimum width=2.35cm,
    minimum height=2.15cm,fill=STBlue!4] (R) at (5.35,0) {};
  \node[font=\fontsize{12pt}{14pt}\selectfont\bfseries,text=STInk] at (1.35,0.35)
    {$B_{\ell-1}^{L}$};
  \node[font=\fontsize{12pt}{14pt}\selectfont\bfseries,text=STInk] at (5.35,0.35)
    {$B_{\ell-1}^{R}$};
  \foreach \x in {0.78,1.35,1.92}
    \node[stDet,minimum size=0.35cm] at (\x,-0.40) {};
  \foreach \x in {4.78,5.35,5.92}
    \node[stDet,minimum size=0.35cm] at (\x,-0.40) {};
  \node[stSmallLabel] at (1.35,-0.78) {ports};
  \node[stSmallLabel] at (5.35,-0.78) {ports};
  \node[stStoch,minimum size=0.72cm] (M) at (3.35,-2.10) {$m_\ell$};

  \draw[stArrow,draw=STBlue,line width=1.3pt]
    (L.east)--node[above,stLabel,text=STBlue] {$w_\ell$} (R.west);
  \draw[stArrow,draw=STCoral,line width=1.3pt]
    (R.south) to[bend left=12]
    node[right,stLabel,text=STCoral] {$w_\ell$} (M.east);
  \draw[stArrow,draw=STTeal,line width=1.3pt]
    (M.west) to[bend left=12]
    node[left,stLabel,text=STTeal] {$w_\ell$} (L.south);

  \node[stDepot,minimum size=0.62cm] (D) at (3.35,1.85) {$r$};
  \draw[stArrowThin,draw=STGold!75!STInk] (D) to[bend right=12] (L.north);
  \draw[stArrowThin,draw=STGold!75!STInk] (L.north) to[bend right=12] (D);
  \draw[stArrowThin,draw=STGold!75!STInk] (D) to[bend left=12] (R.north);
  \draw[stArrowThin,draw=STGold!75!STInk] (R.north) to[bend left=12] (D);
  \node[stSmallLabel,text=STGold!70!STInk] at (3.35,1.22)
    {depot arcs also $w_\ell$};
\end{tikzpicture}}

\newcommand{\StLocalChoice}{%
\begin{tikzpicture}[x=1cm,y=1cm]
  \node[anchor=south,align=center,
    font=\fontsize{8.4pt}{10pt}\selectfont\bfseries,text=STBlue]
    at (1.35,1.90) {KNOWN: \(M\) ABSENT};
  \node[stDet] (L0) at (0.45,0.90) {$L$};
  \node[stDet] (R0) at (2.25,0.90) {$R$};
  \node[stStochInactive] (M0) at (1.35,-0.35) {$M$};
  \draw[stArrow,draw=STBlue,line width=1.35pt] (L0)--node[above,stSmallLabel] {$w$} (R0);
  \node[font=\fontsize{15pt}{17pt}\selectfont\bfseries,text=STInk]
    at (1.35,-1.05) {$w$};

  \draw[STLine] (3.25,-1.35)--(3.25,2.35);
  \node[anchor=south,align=center,
    font=\fontsize{8.4pt}{10pt}\selectfont\bfseries,text=STTeal]
    at (5.00,1.90) {KNOWN: \(M\) ACTIVE};
  \node[stDet] (L1) at (4.10,0.90) {$L$};
  \node[stDet] (R1) at (5.90,0.90) {$R$};
  \node[stStochActive] (M1) at (5.00,-0.35) {$M$};
  \draw[stArrow,draw=STTeal,line width=1.35pt] (R1)--(M1);
  \draw[stArrow,draw=STTeal,line width=1.35pt] (M1)--(L1);
  \node[font=\fontsize{15pt}{17pt}\selectfont\bfseries,text=STInk]
    at (5.00,-1.05) {$2w$};

  \draw[STLine] (6.90,-1.35)--(6.90,2.35);
  \node[anchor=south,align=center,
    font=\fontsize{8.4pt}{10pt}\selectfont\bfseries,text=STCoral]
    at (8.65,1.90) {CAUSAL: PROBE \(M\) LAST};
  \node[stDet] (L2) at (7.75,0.90) {$L$};
  \node[stDet] (R2) at (9.55,0.90) {$R$};
  \node[stStoch] (M2) at (8.65,-0.35) {$M$};
  \draw[stArrow,draw=STBlue,line width=1.25pt] (L2)--(R2);
  \draw[stArrow,draw=STCoral,line width=1.25pt] (R2)--(M2);
  \draw[stArrow,draw=STCoral,line width=1.25pt] (M2)--(L2);
  \node[font=\fontsize{15pt}{17pt}\selectfont\bfseries,text=STInk]
    at (8.65,-1.05) {$\E=2w$};
\end{tikzpicture}}

% --------------------------------------------------------------------------
% Circulation: backbone, fan construction, boosting, and rounding
% --------------------------------------------------------------------------
\newcommand{\StBackboneSample}{%
\begin{tikzpicture}[scale=0.95]
  \foreach \i in {1,...,8}
    \coordinate (b\i) at ({90-45*(\i-1)}:2.55);
  \foreach \i in {1,...,8} {
    \pgfmathtruncatemacro{\next}{mod(\i,8)+1}
    \draw[STBlue!75,line width=1.15pt] (b\i)--(b\next);
  }
  \foreach \i in {1,...,8}
    \node[stDet,minimum size=0.44cm] at (b\i) {};
  \node[stDepot] at (b1) {$r$};
  \node[stStoch,minimum size=0.46cm] (u1) at (-0.25,0.72) {$i$};
  \node[stStoch,minimum size=0.46cm] (u2) at (0.65,-0.72) {$j$};
  \node[stStoch,minimum size=0.46cm] (u3) at (-0.95,-0.48) {$k$};
  \draw[stDashed,draw=STTeal] (b7)--(u1)--(b2);
  \draw[stDashed,draw=STTeal] (b5)--(u2)--(b3);
  \draw[stDashed,draw=STTeal] (b6)--(u3)--(b8);
  \node[stLabel,text=STBlue,anchor=west] at (3.05,1.25)
    {$R=S\cup\{r\}$\\backbone};
  \node[stLabel,text=STTeal,anchor=west] at (3.05,-0.55)
    {$U=V\setminus S$\\outside clients};
\end{tikzpicture}}

\newcommand{\StFanFlow}{%
\begin{tikzpicture}[x=1cm,y=1cm]
  \node[stDet] (s1) at (0,0) {$s_j$};
  \node[stStoch] (u1) at (1.45,0) {$u_1$};
  \node[stStoch] (u2) at (2.75,0) {$u_2$};
  \node[font=\fontsize{13pt}{14pt}\selectfont,text=STMuted] at (3.8,0) {$\cdots$};
  \node[stStoch] (uk) at (4.75,0) {$u_k$};
  \node[stDet] (s2) at (6.25,0) {$s_{j+1}$};
  \draw[STLine,line width=1.3pt] (s1)--(u1)--(u2);
  \draw[STLine,line width=1.3pt] (u2)--(3.43,0);
  \draw[STLine,line width=1.3pt] (4.17,0)--(uk)--(s2);
  \node[stLabel,anchor=south,text=STInkSoft] at (3.10,1.30)
    {segment of an optimal tour on $C=A\cup S$};

  \draw[stArrow,draw=STTeal] (s1.north) to[bend left=24] (u1.north);
  \draw[stArrow,draw=STTeal] (u1.north) to[bend left=24] (s2.north);
  \draw[stArrow,draw=STBlue] (s1.north) to[bend left=42] (u2.north);
  \draw[stArrow,draw=STBlue] (u2.north) to[bend left=42] (s2.north);
  \draw[stArrow,draw=STCoral] (s1.south) to[bend right=28] (uk.south);
  \draw[stArrow,draw=STCoral] (uk.south) to[bend right=28] (s2.south);
  \node[stLabel,anchor=north,text=STInkSoft] at (3.12,-1.35)
    {fan every outside client to the same two backbone endpoints};
\end{tikzpicture}}

\newcommand{\StLongestRun}{%
\begin{tikzpicture}[x=0.92cm,y=1cm]
  \foreach \x in {0,...,11}{
    \pgfmathtruncatemacro{\styleflag}{mod(\x+2,5)}
    \ifnum\styleflag=0
      \node[stDet,minimum size=0.42cm] (n\x) at (\x,0) {};
    \else
      \node[stStoch,minimum size=0.42cm] (n\x) at (\x,0) {};
    \fi
  }
  \foreach \x [evaluate=\x as \next using int(\x+1)] in {0,...,10}
    \draw[STLine,line width=0.75pt] (n\x)--(n\next);
  \draw[decorate,decoration={brace,amplitude=5pt},STCoral,line width=0.9pt]
    (n5.north west)--(n8.north east)
    node[midway,above=0.24cm,text=STCoral,
      font=\fontsize{10pt}{12pt}\selectfont\bfseries] {run length $t$};
  \node[stLabel,anchor=north] at (5.5,-0.50)
    {conditioned on $C$, every vertex is outside with probability at most $\frac12$};
  \node[anchor=south east,text=STInk,
    font=\fontsize{15pt}{17pt}\selectfont\bfseries] at (10.95,1.45)
    {$\Pr[M\ge t\mid C]\le n\,2^{-t}$};
\end{tikzpicture}}

\newcommand{\StBoostedBackbone}{%
\begin{tikzpicture}[x=1cm,y=1cm]
  \foreach \row in {0,1,2}{
    \foreach \x in {0,...,8}{
      \pgfmathtruncatemacro{\isblue}{mod(2*\x+\row,7)}
      \ifnum\isblue<2
        \node[circle,fill=STBlue,minimum size=0.19cm] at (\x*0.75,2.0-\row*0.62) {};
      \else
        \node[circle,draw=STLine,fill=STPaper,minimum size=0.18cm]
          at (\x*0.75,2.0-\row*0.62) {};
      \fi
    }
    \node[anchor=east,stSmallLabel] at (-0.25,2.0-\row*0.62) {$B^{\row}$};
  }
  \draw[decorate,decoration={brace,amplitude=5pt,mirror},STBlue,line width=0.8pt]
    (-0.25,0.48)--(6.35,0.48)
    node[midway,below=0.23cm,stLabel,text=STBlue]
    {$S=\bigcup_{j=1}^{\lambda}B^j$,\quad $\lambda=\Theta(\log n)$};
  \draw[stArrow,draw=STTeal] (6.90,1.25)--(8.10,1.25);
  \node[anchor=west,align=left,text=STInk,
    font=\fontsize{10pt}{12.5pt}\selectfont] at (8.35,1.25)
    {\textbf{Outside probabilities shrink}\\[1mm]
     $p_i<1/\lambda$};
\end{tikzpicture}}

\newcommand{\StChordCoupling}{%
\begin{tikzpicture}[x=1cm,y=1cm]
  \node[stDet] (a1) at (0,1.55) {$a_1$};
  \node[stDet] (a2) at (0,0.35) {$a_2$};
  \node[stDet] (a3) at (0,-0.85) {$a_3$};
  \node[stStoch,minimum size=0.64cm] (i) at (3.2,0.35) {$i$};
  \node[stDet] (b1) at (6.4,1.55) {$b_1$};
  \node[stDet] (b2) at (6.4,0.35) {$b_2$};
  \node[stDet] (b3) at (6.4,-0.85) {$b_3$};
  \draw[stArrowThin,draw=STBlue] (a1)--(i);
  \draw[stArrowThin,draw=STBlue] (a2)--(i);
  \draw[stArrowThin,draw=STBlue] (a3)--(i);
  \draw[stArrowThin,draw=STTeal] (i)--(b1);
  \draw[stArrowThin,draw=STTeal] (i)--(b2);
  \draw[stArrowThin,draw=STTeal] (i)--(b3);
  \node[stSmallLabel,text=STBlue,anchor=south] at (1.25,1.72)
    {incoming marginal $z_{ai}$};
  \node[stSmallLabel,text=STTeal,anchor=south] at (5.10,1.72)
    {outgoing marginal $z_{ib}$};
  \draw[stArrow,draw=STCoral,line width=1.25pt]
    (0.10,-1.65) to[bend right=15]
    node[below,stLabel,text=STCoral] {couple into $x_{i,ab}$}
    (6.30,-1.65);
  \node[anchor=west,align=left,text=STInk,
    font=\fontsize{10pt}{12pt}\selectfont] at (7.05,0.70)
    {$\displaystyle\sum_{a,b}x_{i,ab}=p_i$\\[2mm]
     round to one pair\\
     $\widehat x_{i,a_i b_i}=p_i$};
\end{tikzpicture}}

\newcommand{\StDyadicIntervals}{%
\begin{tikzpicture}[x=1cm,y=1cm]
  \foreach \x in {0,...,7}{
    \node[stDet,minimum size=0.41cm] (p\x) at (\x*1.10,0) {$s_{\x}$};
  }
  \foreach \x [evaluate=\x as \next using int(\x+1)] in {0,...,6}
    \draw[STLine] (p\x)--(p\next);

  \draw[STTeal,line width=2pt] (-0.28,0.72)--(8.0,0.72);
  \node[anchor=west,stSmallLabel,text=STTeal] at (8.15,0.72) {level 0};
  \foreach \a/\b in {-0.28/4.10,4.10/8.0}
    \draw[STBlue,line width=2pt] (\a,1.30)--(\b,1.30);
  \node[anchor=west,stSmallLabel,text=STBlue] at (8.15,1.30) {level 1};
  \foreach \a/\b in {-0.28/1.93,1.93/4.10,4.10/6.28,6.28/8.0}
    \draw[STGold,line width=2pt] (\a,1.88)--(\b,1.88);
  \node[anchor=west,stSmallLabel,text=STGold!70!STInk] at (8.15,1.88) {level 2};

  \draw[decorate,decoration={brace,amplitude=5pt,mirror},STCoral,line width=0.9pt]
    (p2.south west)--(p6.south east)
    node[midway,below=0.25cm,stLabel,text=STCoral]
      {any interval = $O(\log n)$ dyadic pieces};
  \node[anchor=west,align=left,text=STInkSoft,
    font=\fontsize{9.2pt}{11.5pt}\selectfont] at (0,2.65)
    {Preserve a dyadic balance equation while it touches
     at least $L=8(\lceil\log_2 m\rceil+1)$ live variables.};
\end{tikzpicture}}

\newcommand{\StRoundedChords}{%
\begin{tikzpicture}[scale=0.86]
  \foreach \i in {1,...,8}
    \coordinate (q\i) at ({90-45*(\i-1)}:2.7);
  \foreach \i in {1,...,8} {
    \pgfmathtruncatemacro{\next}{mod(\i,8)+1}
    \draw[STBlue!75,line width=1.2pt,-{Stealth[length=1.5mm]}] (q\i)--(q\next);
    \node[stDet,minimum size=0.40cm] at (q\i) {};
  }
  \draw[stArrow,draw=STTeal,line width=1.1pt] (q7) to[bend left=20] (q2);
  \draw[stArrow,draw=STCoral,line width=1.1pt] (q8) to[bend left=12] (q4);
  \draw[stArrow,draw=STGold!75!STInk,line width=1.1pt] (q5) to[bend left=18] (q1);
  \draw[STCoral,line width=2.7pt] (q8)--(q1)--(q2)--(q3);
  \node[stLabel,text=STCoral,anchor=west] at (3.15,1.10)
    {interval $I$};
  \node[stLabel,text=STInk,anchor=west,align=left] at (3.15,-0.35)
    {$|\Delta(I)|=O(\log n)$\\after boosting + rounding};
\end{tikzpicture}}

\newcommand{\StOfflineWalk}{%
\begin{tikzpicture}[scale=0.88]
  \foreach \i in {1,...,7}
    \coordinate (o\i) at ({90-360/7*(\i-1)}:2.45);
  \foreach \i in {1,...,7} {
    \pgfmathtruncatemacro{\next}{mod(\i,7)+1}
    \draw[STBlue,line width=2.0pt,-{Stealth[length=1.8mm]}] (o\i)--(o\next);
    \node[stDet,minimum size=0.43cm] at (o\i) {};
  }
  \node[stStochActive,minimum size=0.46cm] (i1) at (0.15,0.55) {$i$};
  \node[stStochActive,minimum size=0.46cm] (i2) at (-0.75,-0.62) {$j$};
  \draw[stArrow,draw=STTeal,line width=1.35pt] (o6)--(i1);
  \draw[stArrow,draw=STTeal,line width=1.35pt] (i1)--(o2);
  \draw[stArrow,draw=STCoral,line width=1.35pt] (o5)--(i2);
  \draw[stArrow,draw=STCoral,line width=1.35pt] (i2)--(o3);
  \node[anchor=west,align=left,text=STInk,
    font=\fontsize{10pt}{12.5pt}\selectfont] at (3.20,0.75)
    {\textbf{Hoffman circulation}\\[1mm]
     active chord: $1$--$4$ copies\\
     backbone / auxiliary arc:\\
     $O(\log n)$ copies};
  \node[anchor=west,align=left,text=STCoral,
    font=\fontsize{9.5pt}{11.5pt}\selectfont\bfseries] at (3.20,-1.30)
    {chosen after seeing all of $A$};
\end{tikzpicture}}

% --------------------------------------------------------------------------
% Dynamic programming: pairwise cost, inflation, blocks, and suffix memory
% --------------------------------------------------------------------------
\newcommand{\StPairwiseShortcut}{%
\begin{tikzpicture}[x=1cm,y=1cm]
  \node[stDepot] (r) at (0,0) {$r$};
  \node[stDet] (one) at (1.25,0) {$v_1$};
  \node[stStochActive] (vi) at (2.50,0) {$v_i$};
  \node[stStochInactive] (a) at (3.65,0) {$a$};
  \node[stStochInactive] (b) at (4.80,0) {$b$};
  \node[stStochInactive] (c) at (5.95,0) {$c$};
  \node[stStochActive] (vj) at (7.10,0) {$v_j$};
  \node[stDet] (d) at (8.35,0) {$d$};
  \node[stDepot] (rb) at (9.60,0) {$r$};
  \foreach \a/\b in {r/one,one/vi,vi/a,a/b,b/c,c/vj,vj/d,d/rb}
    \draw[STLine,line width=0.75pt] (\a)--(\b);
  \draw[stArrow,draw=STCoral,line width=1.5pt]
    (vi.north) to[bend left=27]
    node[above,align=center,text=STCoral,
      font=\fontsize{10pt}{12pt}\selectfont\bfseries]
    {$p_i p_j\,Q(i,j)\,d_{ij}$\\[-1mm]
     $Q(i,j)=\prod_{i<k<j}(1-p_k)$}
    (vj.north);
  \node[stSmallLabel,anchor=north] at (4.8,-0.42)
    {all intermediate clients inactive};
\end{tikzpicture}}

\newcommand{\StActivationInflation}{%
\begin{tikzpicture}[x=1cm,y=1cm]
  \node[anchor=west,text=STInk,font=\fontsize{11pt}{13pt}\selectfont\bfseries]
    at (0,2.35) {Original activations};
  \foreach \x in {0,...,7}{
    \pgfmathtruncatemacro{\flag}{mod(3*\x+1,5)}
    \ifnum\flag<2
      \node[stStochActive,minimum size=0.38cm] at (0.4+\x*0.86,1.45) {};
    \else
      \node[stStochInactive,minimum size=0.38cm] at (0.4+\x*0.86,1.45) {};
    \fi
  }
  \draw[stArrow,draw=STBlue] (7.60,1.45)--(8.70,1.45);
  \node[anchor=west,align=left,text=STBlue,
    font=\fontsize{9.3pt}{11.5pt}\selectfont] at (8.95,1.45)
    {$\widetilde p_v=1-(1-p_v)^{1+\alpha}$\\
     $\widetilde Q=Q^{1+\alpha}$};

  \node[anchor=west,text=STInk,font=\fontsize{11pt}{13pt}\selectfont\bfseries]
    at (0,0.42) {After an independent extra round};
  \foreach \x in {0,...,7}{
    \pgfmathtruncatemacro{\flag}{mod(3*\x+1,5)}
    \ifnum\flag<4
      \node[stStochActive,minimum size=0.38cm] at (0.4+\x*0.86,-0.48) {};
    \else
      \node[stStochInactive,minimum size=0.38cm] at (0.4+\x*0.86,-0.48) {};
    \fi
  }
  \draw[decorate,decoration={brace,amplitude=5pt,mirror},STCoral,line width=0.8pt]
    (2.0,-1.02)--(5.65,-1.02)
    node[midway,below=0.22cm,stLabel,text=STCoral]
    {a long empty gap is damped by $Q^\alpha$};
\end{tikzpicture}}

\newcommand{\StBlockStrip}{%
\begin{tikzpicture}[x=1cm,y=1cm]
  \node[stStochInactive] (l1) at (0.30,0) {};
  \node[stStochInactive] (l2) at (0.95,0) {};
  \node[stStochInactive] (l3) at (1.60,0) {};
  \node[stStochActive,minimum size=0.55cm] (h1) at (2.75,0) {$H$};
  \node[stStochInactive] (l4) at (3.85,0) {};
  \node[stStochInactive] (l5) at (4.50,0) {};
  \node[stStochInactive] (l6) at (5.15,0) {};
  \node[stStochInactive] (l7) at (5.80,0) {};
  \node[stStochActive,minimum size=0.55cm] (h2) at (6.95,0) {$H$};
  \node[stStochInactive] (l8) at (8.05,0) {};
  \node[stStochInactive] (l9) at (8.70,0) {};

  \draw[STBlue,line width=1.15pt,rounded corners=0.65pt]
    (0.0,-0.55) rectangle (1.90,0.55);
  \draw[STGold,line width=1.15pt,rounded corners=0.65pt]
    (2.42,-0.55) rectangle (3.08,0.55);
  \draw[STTeal,line width=1.15pt,rounded corners=0.65pt]
    (3.55,-0.55) rectangle (6.10,0.55);
  \draw[STGold,line width=1.15pt,rounded corners=0.65pt]
    (6.62,-0.55) rectangle (7.28,0.55);
  \draw[STBlue,line width=1.15pt,rounded corners=0.65pt]
    (7.75,-0.55) rectangle (9.00,0.55);
  \node[stSmallLabel,text=STBlue] at (0.95,-0.90) {light block};
  \node[stSmallLabel,text=STGold!75!STInk] at (2.75,-0.90) {heavy singleton};
  \node[stSmallLabel,text=STTeal] at (4.82,-0.90)
    {$\delta/2\le\mu(B)\le\delta$};
  \node[stSmallLabel,text=STBlue] at (8.38,-0.90)
    {first light block\\may be smaller};
\end{tikzpicture}}

\newcommand{\StDPSuffix}{%
\begin{tikzpicture}[x=1cm,y=1cm]
  \fill[STLine!45,rounded corners=0.65pt] (-0.15,-0.48) rectangle (3.35,0.58);
  \node[text=STMuted,font=\fontsize{10pt}{12pt}\selectfont\bfseries]
    at (1.60,0.05) {placed and forgotten};
  \draw[STLine,dashed,line width=1pt] (3.65,-1.05)--(3.65,1.28);
  \node[anchor=south,text=STInkSoft,font=\fontsize{8.5pt}{10pt}\selectfont]
    at (3.65,1.30) {$\prod q\le\eta$};
  \node[draw=STBlue,fill=STBlue!9,rounded corners=0.65pt,
    minimum width=1.05cm,minimum height=0.76cm,
    font=\fontsize{9pt}{10pt}\selectfont\bfseries,text=STInk]
    (Ba) at (4.35,0.05) {$B_a$};
  \node[draw=STTeal,fill=STTeal!9,rounded corners=0.65pt,
    minimum width=1.05cm,minimum height=0.76cm,
    font=\fontsize{9pt}{10pt}\selectfont\bfseries,text=STInk]
    (Bap) at (5.70,0.05) {$B_{a+1}$};
  \node[draw=STMuted,fill=STMuted!9,rounded corners=0.65pt,
    minimum width=1.05cm,minimum height=0.76cm,
    font=\fontsize{9pt}{10pt}\selectfont\bfseries,text=STInk]
    (Bdots) at (7.05,0.05) {$\cdots$};
  \node[draw=STGold,fill=STGold!9,rounded corners=0.65pt,
    minimum width=1.05cm,minimum height=0.76cm,
    font=\fontsize{9pt}{10pt}\selectfont\bfseries,text=STInk]
    (Bt) at (8.40,0.05) {$B_t$};
  \draw[stArrow,draw=STCoral] (9.05,0.05)--(10.15,0.05);
  \node[draw=STCoral,fill=STCoral!10,rounded corners=0.65pt,
    minimum width=1.06cm,minimum height=0.76cm,
    font=\fontsize{9pt}{10pt}\selectfont\bfseries] (B) at (10.85,0.05) {$B$};
  \foreach \source in {Ba,Bap,Bt}
    \draw[stArrowThin,draw=STBlue!75] (\source.north)
      to[bend left=20] (B.north);
  \node[anchor=north,align=center,text=STInkSoft,
    font=\fontsize{9pt}{11pt}\selectfont] at (7.55,-0.83)
    {state $(U,\Omega)$ remembers only the relevant suffix $\Omega$};
\end{tikzpicture}}

\setbeameroption{hide notes}

\begin{document}

% ==========================================================================
% Title
% ==========================================================================
\begin{frame}[plain]
  \begin{tikzpicture}[remember picture,overlay]
    \fill[STPaper] (current page.south west) rectangle (current page.north east);
    \draw[STLine,line width=0.45pt]
      ([xshift=0.95cm,yshift=-0.82cm]current page.north west)
      --
      ([xshift=-0.88cm,yshift=-0.82cm]current page.north east);
    \draw[STBlue,line width=1.35pt]
      ([xshift=0.95cm,yshift=-0.82cm]current page.north west)
      --
      ([xshift=3.35cm,yshift=-0.82cm]current page.north west);

    \node[
      anchor=west,
      text=STBlue,
      font=\fontsize{9pt}{11pt}\selectfont\bfseries
    ] at ([xshift=0.98cm,yshift=2.42cm]current page.west)
      {MASTER THESIS \enspace·\enspace JULY 2026};
    \node[
      anchor=east,
      text=STMuted,
      font=\fontsize{9.4pt}{11.5pt}\selectfont
    ] at ([xshift=-0.90cm,yshift=2.42cm]current page.east)
      {\(\optpost\leq\optadapt\leq\optprior\)};
    \node[
      anchor=west,
      text=STInk,
      text width=11.35cm,
      align=left,
      font=\fontsize{20pt}{23.8pt}\selectfont\bfseries
    ] at ([xshift=0.98cm,yshift=0.34cm]current page.west)
      {Information Gaps and Algorithms\\
       for the Stochastic Traveling\\
       Salesperson Problem};
    \draw[STLine,line width=0.45pt]
      ([xshift=0.98cm,yshift=-1.10cm]current page.west)
      --
      ([xshift=-0.90cm,yshift=-1.10cm]current page.east);
    \node[
      anchor=west,
      text=STInk,
      font=\fontsize{13.2pt}{16pt}\selectfont\bfseries
    ] at ([xshift=0.98cm,yshift=-1.54cm]current page.west)
      {Teng Liu};
    \node[
      anchor=west,
      text=STInkSoft,
      font=\fontsize{10.8pt}{13pt}\selectfont
    ] at ([xshift=0.98cm,yshift=-1.98cm]current page.west)
      {Department of Computer Science \enspace·\enspace ETH Zürich};
    \node[
      anchor=west,
      text=STMuted,
      font=\fontsize{9.8pt}{12pt}\selectfont
    ] at ([xshift=0.98cm,yshift=-2.56cm]current page.west)
      {Advisors: Dr.\ Sharat Ibrahimpur and Prof.\ Dr.\ Vera Traub};
  \end{tikzpicture}
\end{frame}

% ==========================================================================
% Outline — outside the requested technical timing
% ==========================================================================
\begin{frame}{Outline}
  \vfill
  \tableofcontents[hideallsubsections]
  \vfill
\end{frame}

% ==========================================================================
% Background — 3 minutes
% ==========================================================================
\section{Problem and benchmarks}

% Timing: ~1:00
\begin{frame}{Uncertainty is simple; ordering is not}
  \begin{columns}[T,onlytextwidth]
    \begin{column}{0.35\textwidth}
      \vspace{0.00cm}
      Each client \(v\) is independently active with probability \(p_v\).

      \vspace{0.24cm}
      A tour starts and ends at depot \(r\), visits every active client, and
      pays directed metric distance.

      \vspace{0.22cm}
      {\small
      \stresult[STTeal]{The hard decision is \emph{when the visit order may
      depend on the realization}.}}
    \end{column}
    \begin{column}{0.62\textwidth}
      \centering
      \resizebox{0.99\linewidth}{!}{\StMasterOrder}
    \end{column}
  \end{columns}
  \vspace{-0.02cm}
\end{frame}

% Timing: ~1:10
\begin{frame}{Information arrives at three different times}
  \centering
  \resizebox{0.74\linewidth}{!}{\StBenchmarkLadder}

  \vspace{0.02cm}
  \[
    \optpost \;\le\; \optadapt \;\le\; \optprior
  \]

  \vspace{-0.30cm}
  \small
  \stpill[STBlue]{adaptivity gap}
  \(\displaystyle\sup\frac{\optprior}{\optadapt}\)
  \hspace{0.75cm}
  \stpill[STTeal]{clairvoyance gap}
  \(\displaystyle\sup\frac{\optadapt}{\optpost}\)
\end{frame}

% Timing: ~0:50
\begin{frame}{The thesis separates three kinds of difficulty}
  \vspace{0.18cm}
  \begin{tikzpicture}[x=1cm,y=1cm]
    \draw[STLine,line width=1.5pt] (0.75,2.20)--(10.95,2.20);
    \foreach \x/\col/\label in {
      1.20/STBlue/1,
      5.80/STTeal/2,
      10.45/STGold/3
    }{
      \node[circle,fill=\col,text=STWhite,minimum size=0.58cm,
        font=\fontsize{11pt}{12pt}\selectfont\bfseries] at (\x,2.20) {\label};
    }
    \node[anchor=north west,text width=3.25cm,align=left] at (0.75,1.70) {
      \textbf{Witness the gaps}\par
      \vspace{0.16cm}
      \color{STInkSoft}\small
      Symmetric: \(5/4\), \(11/10\)\\
      Asymmetric: \(4/3\)
    };
    \node[anchor=north,text width=3.45cm,align=center] at (5.80,1.70) {
      \textbf{Build low-cost structure}\par
      \vspace{0.16cm}
      \color{STInkSoft}\small
      Backbone circulation,\\endpoint rounding, offline walk
    };
    \node[anchor=north east,text width=3.15cm,align=right] at (10.95,1.70) {
      \textbf{Optimize a fixed order}\par
      \vspace{0.16cm}
      \color{STInkSoft}\small
      Factor \(1+\varepsilon\)\\
      Time \(c_\varepsilon^n\poly(n)\)
    };
  \end{tikzpicture}

  \vfill
  \centering
  \large
  The common theme is not merely uncertainty—it is the value and the
  \emph{timing} of information.
\end{frame}

% ==========================================================================
% Constructions — 13 minutes
% ==========================================================================
\section{Gap constructions}
\stsectionframe{STBlue}{Part I}{Three gap constructions}{Local switching,
global restructuring,\\and recursive reversal.}

% Timing: ~1:05
\begin{frame}{Subdividing a cycle makes observations useful}
  \begin{columns}[T,onlytextwidth]
    \begin{column}{0.55\textwidth}
      \centering
      \resizebox{!}{5.00cm}{\StCycleClique}
    \end{column}
    \begin{column}{0.42\textwidth}
      \vspace{0.42cm}
      Start from \(K_n\), then subdivide the edges
      \(v_i v_{i+1}\) of one Hamiltonian cycle.

      \vspace{0.38cm}
      \begin{sttightitemize}
        \item original vertices: deterministic;
        \item subdividers \(v_i'\): \(p=\tfrac12\);
        \item distances: shortest-path metric.
      \end{sttightitemize}

      \vspace{0.30cm}
      \stpill[STBlue]{symmetric metric}
    \end{column}
  \end{columns}
\end{frame}

% Timing: ~1:10
\begin{frame}{Every fixed order pays a local tax}
  \begin{columns}[T,onlytextwidth]
    \begin{column}{0.58\textwidth}
      \centering
      \vspace{0.42cm}
      \resizebox{0.94\linewidth}{!}{\StSegmentTax}
    \end{column}
    \begin{column}{0.39\textwidth}
      A fixed tour breaks into \(n\) segments. If segment \(j\) contains
      \(k_j\) subdividers, then
      \[
        \E[\cost(T)]
        \ge n+\frac78\sum_{j=1}^{n}k_j.
      \]

      Because \(\sum_j k_j=n\),
      \[
        \boxed{\optprior=\frac{15}{8}n.}
      \]
    \end{column}
  \end{columns}
\end{frame}

% Timing: ~1:05
\begin{frame}{Adaptation peels from the cheaper end}
  \begin{columns}[T,onlytextwidth]
    \begin{column}{0.57\textwidth}
      \centering
      \vspace{0.62cm}
      \resizebox{0.95\linewidth}{!}{\StAdaptivePeel}
    \end{column}
    \begin{column}{0.40\textwidth}
      Maintain the uncalled part of the subdivided path as one interval.

      \vspace{0.12cm}
      \begin{sttightitemize}
        \item active: cross it, cost \(2\);
        \item inactive: switch endpoints, usually cost \(1\).
      \end{sttightitemize}

      \vspace{0.08cm}
      With \(a\) active internal subdividers,
      {\small
      \[
        \begin{aligned}
          \cost(\pi,A)&\le n+a+7,\\[-1mm]
          \optadapt&\le \frac32n+6.
        \end{aligned}
      \]}
    \end{column}
  \end{columns}
\end{frame}

% Timing: ~0:50
\begin{frame}{Local switching creates a \(5/4\) gap}
  \vspace{0.70cm}
  \centering
  \steyebrow{fixed order versus adaptive policy}

  \vspace{0.30cm}
  {\fontsize{29pt}{34pt}\selectfont\bfseries
    \(\displaystyle
      \frac{\optprior}{\optadapt}
      \;\ge\;
      \frac{\frac{15}{8}n}{\frac32n+6}
      \;\longrightarrow\;
      \color{STBlue}\frac54
    \)}

  \vspace{0.64cm}
  \centering
  The gain comes from changing direction after an inactive call.
\end{frame}

% Timing: ~1:10
\begin{frame}{One stochastic vertex forces global restructuring}
  \begin{columns}[T,onlytextwidth]
    \begin{column}{0.61\textwidth}
      \centering
      \resizebox{0.96\linewidth}{!}{\StThreeChain}
    \end{column}
    \begin{column}{0.36\textwidth}
      Put three unit-edge chains of length \(h=k-1\) behind seven inter-chain
      edges of length \(2h\).

      \vspace{0.38cm}
      \begin{sttightitemize}
        \item depot: \(a_0\);
        \item all chain vertices deterministic;
        \item only \(s\) is stochastic, with \(p_s=\tfrac12\).
      \end{sttightitemize}

      \vspace{0.28cm}
      \stpill[STTeal]{symmetric metric}
    \end{column}
  \end{columns}
\end{frame}

% Timing: ~1:00
\begin{frame}{The realizations prefer disjoint crossing skeletons}
  \centering
  \resizebox{0.79\linewidth}{!}{\StSkeletonComparison}

  \vspace{0.02cm}
  {\small
  \stresult[STCoral]{Removing \(s\) reverses the preferred order of
  entire chains.}}
\end{frame}

% Timing: ~1:15
\begin{frame}{Adaptation commits before learning the skeleton}
  \begin{columns}[T,onlytextwidth]
    \begin{column}{0.47\textwidth}
      Near-optimal tours are forced:
      \[
      \begin{aligned}
        C_0<11h &\Longrightarrow \text{inactive skeleton }E_D,\\
        C_1<13h &\Longrightarrow \text{active skeleton }E_{\mathrm{all}}.
      \end{aligned}
      \]
      The two crossing sets are disjoint.
    \end{column}
    \begin{column}{0.48\textwidth}
      Before calling \(s\), both realizations have the same called prefix.

      \vspace{0.30cm}
      That prefix cannot cross between chains in either forced skeleton.
      Reaching \(s\) then transits through an uncalled chain with no crossing
      left to return and serve it.

      \vspace{0.28cm}
      \stpill[STCoral]{contradiction}
    \end{column}
  \end{columns}

  \vfill
  \centering
  Therefore both \(C_0<11h\) and \(C_1<13h\) are impossible.
\end{frame}

% Timing: ~0:50
\begin{frame}{Clairvoyance saves one expected crossing}
  \vspace{0.22cm}
  \begin{columns}[c,onlytextwidth]
    \begin{column}{0.46\textwidth}
      \centering
      \steyebrow{a posteriori}
      \[
        \optpost=\tfrac12(9h)+\tfrac12(11h)=10h
      \]
    \end{column}
    \begin{column}{0.46\textwidth}
      \centering
      \steyebrow{adaptive}
      \[
        \optadapt=11h
      \]
    \end{column}
  \end{columns}

  \vspace{0.10cm}
  \centering
  {\fontsize{31pt}{35pt}\selectfont\bfseries\color{STTeal}
    \(\displaystyle \frac{\optadapt}{\optpost}=\frac{11}{10}\)}

  \vspace{0.28cm}
  Here \(\optprior=\optadapt\): the entire gap is the value of knowing \(s\)
  before choosing the chain order.
\end{frame}

% Timing: ~1:05
\begin{frame}{A recursive block repeats the order reversal}
  \begin{columns}[T,onlytextwidth]
    \begin{column}{0.58\textwidth}
      \centering
      \resizebox{0.79\linewidth}{!}{\StRecursiveBlock}
    \end{column}
    \begin{column}{0.39\textwidth}
      Build \(B_\ell\) from two copies of \(B_{\ell-1}\) and midpoint \(m_\ell\).

      \[
        \begin{gathered}
          L\longrightarrow R\longrightarrow m_\ell\longrightarrow L,\\[-1mm]
          w_\ell=2^{\ell-1}.
        \end{gathered}
      \]

      Permanent ports are deterministic; every midpoint is active with
      probability \(1/2\).

      \vspace{0.28cm}
      \stpill[STCoral]{directed metric}
    \end{column}
  \end{columns}
\end{frame}

% Timing: ~1:00
\begin{frame}{An active midpoint reverses the local order}
  \centering
  \resizebox{0.88\linewidth}{!}{\StLocalChoice}

  \vspace{0.02cm}
  \begin{columns}[T,onlytextwidth]
    \begin{column}{0.47\textwidth}
      \centering
      Clairvoyance pays
      \(\tfrac12w+\tfrac12(2w)=\tfrac32w\).
    \end{column}
    \begin{column}{0.47\textwidth}
      \centering
      A causal policy pays
      \(\tfrac12w+\tfrac12(3w)=2w\).
    \end{column}
  \end{columns}
\end{frame}

% Timing: ~1:15
\begin{frame}{Interrupted services obey clean recurrences}
  \begin{columns}[T,onlytextwidth]
    \begin{column}{0.48\textwidth}
      \steyebrow{known realization}
      \[
        \begin{aligned}
        P_\ell(A)
          &=P_{\ell-1}(A_L)+P_{\ell-1}(A_R)\\[-1mm]
          &\quad+w_\ell\bigl(1+\1_{\{m_\ell\in A\}}\bigr).
        \end{aligned}
      \]
      Hence
      \[
        \E[P_\ell]=3\ell\,2^{\ell-2}.
      \]
    \end{column}
    \begin{column}{0.48\textwidth}
      \steyebrow{adaptive open service}
      \[
        Q_\ell=2Q_{\ell-1}+2w_\ell,
        \qquad Q_0=0.
      \]
      Hence
      \[
        Q_\ell=\ell\,2^\ell.
      \]
    \end{column}
  \end{columns}

  \vfill
  \small
  Every additional excursion into a child is charged to a top-level entry
  arc, so interleaving cannot beat these recurrences.
\end{frame}

% Timing: ~1:00
\begin{frame}{The penalty accumulates once per depth}
  \vspace{0.15cm}
  \[
    \optpost(I_\ell)=(3\ell+4)2^{\ell-2},
    \qquad
    \optadapt(I_\ell)=(\ell+1)2^\ell.
  \]

  \vspace{0.24cm}
  \centering
  {\fontsize{30pt}{34pt}\selectfont\bfseries\color{STCoral}
    \(\displaystyle
      \frac{\optadapt(I_\ell)}{\optpost(I_\ell)}
      =\frac{4(\ell+1)}{3\ell+4}
      \longrightarrow \frac43
    \)}

  \vspace{0.24cm}
  \begin{tikzpicture}
    \node[anchor=east,text=STBlue,
      font=\fontsize{11pt}{13pt}\selectfont\bfseries] at (-0.35,0)
      {cycle-subdivided clique};
    \node[anchor=west,text=STInkSoft,
      font=\fontsize{11pt}{13pt}\selectfont] at (0,0)
      {local switching};
    \node[anchor=east,text=STTeal,
      font=\fontsize{10.5pt}{12pt}\selectfont\bfseries] at (-0.35,-0.52)
      {three chains};
    \node[anchor=west,text=STInkSoft,
      font=\fontsize{10.5pt}{12pt}\selectfont] at (0,-0.52)
      {global restructuring};
    \node[anchor=east,text=STCoral,
      font=\fontsize{10.5pt}{12pt}\selectfont\bfseries] at (-0.35,-1.04)
      {recursive blocks};
    \node[anchor=west,text=STInkSoft,
      font=\fontsize{10.5pt}{12pt}\selectfont] at (0,-1.04)
      {the same causal penalty at every scale};
  \end{tikzpicture}
\end{frame}

% ==========================================================================
% Circulation — 15 minutes
% ==========================================================================
\section{Backbone circulation and rounding}
\stsectionframe{STTeal}{Part II}{Directed structure versus\\causal
execution}{A backbone yields a fractional circulation, endpoint rounding,
and an offline walk; causality remains the obstacle.}

% Timing: ~1:00
\begin{frame}{Directed attachments need two endpoints}
  \begin{columns}[T,onlytextwidth]
    \begin{column}{0.58\textwidth}
      \centering
      \resizebox{0.90\linewidth}{!}{\StBackboneSample}
    \end{column}
    \begin{column}{0.39\textwidth}
      Sample \(S\) independently with
      \(\Pp(v\in S)=p_v\).
      \[
        R=S\cup\{r\},\qquad U=V\setminus S.
      \]

      In a symmetric metric, an outside client can detour to one nearby
      backbone point.

      \vspace{0.28cm}
      In a directed metric, entering and leaving may require
      \emph{different} endpoints.
    \end{column}
  \end{columns}
\end{frame}

% Timing: ~1:05
\begin{frame}{The LP is a fractional directed detour system}
  \begin{columns}[T,onlytextwidth]
    \begin{column}{0.59\textwidth}
      {\fontsize{12.5pt}{15pt}\selectfont
      \[
      \begin{aligned}
        \min\quad &\sum_{u,v} d(u,v)z_{uv}\\
        \text{s.t.}\quad
          &z(\delta^+(w))=z(\delta^-(w))
            &&\forall w,\\
          &z(\delta^+(i))=z(\delta^-(i))=p_i
            &&\forall i\in U,\\
          &z_{ij}=0
            &&\forall i,j\in U,\\
          &z_{uv}\ge0
            &&\forall u,v.
      \end{aligned}
      \tag{\(\mathrm{CIRC}(S)\)}
      \]}
    \end{column}
    \begin{column}{0.37\textwidth}
      \vspace{0.40cm}
      No outside-to-outside arc is allowed.

      \vspace{0.38cm}
      Every \(i\in U\) receives \(p_i\) units from the backbone and sends
      \(p_i\) units back.

      \vspace{0.28cm}
      {\small
      \stresult[STTeal]{The LP preserves direction before rounding.}}
    \end{column}
  \end{columns}
\end{frame}

% Timing: ~1:05
\begin{frame}{A realized tour fans into a feasible circulation}
  \begin{columns}[T,onlytextwidth]
    \begin{column}{0.58\textwidth}
      \centering
      \vspace{0.18cm}
      \resizebox{0.96\linewidth}{!}{\StFanFlow}
    \end{column}
    \begin{column}{0.39\textwidth}
      Take an independent activation set \(A\), put \(C=A\cup S\), and cut an
      optimal tour \(T_C\) at backbone vertices.

      \vspace{0.20cm}
      Let \(M(A,S)\) be the longest run of outside vertices between two
      consecutive backbone vertices.

      \[
        \cost(z^{A,S})
        \le M(A,S)\,\OPT(C).
      \]

    \end{column}
  \end{columns}
\end{frame}

% Timing: ~1:05
\begin{frame}{The longest run controls expected cost}
  \centering
  \resizebox{0.80\linewidth}{!}{\StLongestRun}

  \vspace{0.08cm}
  {\fontsize{12pt}{14.5pt}\selectfont
  \[
    \E[M\mid C]=O(\log(n+1))
    \quad\Longrightarrow\quad
    \E_S[\OPT(S)+c(S)]
    \le O(\log(n+1))\optpost.
  \]}
\end{frame}

% Timing: ~1:00
\begin{frame}{Boosting buys regularity for rounding}
  \centering
  \resizebox{0.94\linewidth}{!}{\StBoostedBackbone}

  \vspace{0.24cm}
  With \(q_v=\min\{1,\lambda p_v\}\) and
  \(\lambda=\Theta(\log(n+1))\), some backbone satisfies
  \[
  \begin{aligned}
    \cost(T_R)+\cost(z)
      &\le O(\log(n+1))\optpost,\\
    z(\delta^\pm(s))
      &\le O(\log(n+1)) &&(s\in R),\\
    p_i&<1/\lambda &&(i\notin S).
  \end{aligned}
  \]
\end{frame}

% Timing: ~1:00
\begin{frame}{Couple each client into a chord distribution}
  \centering
  \resizebox{0.58\linewidth}{!}{\StChordCoupling}

  \vspace{-0.02cm}
  {\fontsize{11.8pt}{14pt}\selectfont
  \[
    \sum_b x_{i,ab}=z_{ai},
    \qquad
    \sum_a x_{i,ab}=z_{ib},
    \qquad
    \sum_{a,b}x_{i,ab}=p_i.
  \]}
  The coupling preserves the circulation cost exactly.
\end{frame}

% Timing: ~1:05
\begin{frame}{Preserve a laminar family of balances}
  \centering
  \resizebox{0.94\linewidth}{!}{\StDyadicIntervals}

  \vspace{0.10cm}
  For a cyclic backbone interval \(I\), let
  \[
    \Delta_{x,y}(I)
    =\text{weighted chord flow leaving \(I\)}
     -\text{weighted chord flow entering \(I\)}.
  \]
  Initially every interval has \(\Delta_{x,y}(I)=0\). The rounding preserves
  dyadic interval equations only while they touch many live variables.
\end{frame}

% Timing: ~1:10
\begin{frame}{A dimension count leaves room to move}
  Let \(F=\{x_{i,ab}:0<x_{i,ab}<p_i\}\) be the live variables.

  \vspace{0.22cm}
  \begin{columns}[T,onlytextwidth]
    \begin{column}{0.46\textwidth}
      Every live outside equation contains at least two live variables:
      \[
        c_1\le\frac{|F|}{2}.
      \]
    \end{column}
    \begin{column}{0.46\textwidth}
      A live variable touches at most \(2(h+1)\) dyadic intervals; with
      \(L=8(h+1)\):
      \[
        c_2L\le2(h+1)|F|
        \quad\Rightarrow\quad
        c_2\le\frac{|F|}{4}.
      \]
    \end{column}
  \end{columns}

  \vspace{0.20cm}
  \[
    c_1+c_2<|F|.
  \]
  So a nonzero feasible direction exists. Move until a variable hits
  \(0\) or \(p_i\), choosing the non-increasing-cost endpoint.
\end{frame}

% Timing: ~0:55
\begin{frame}{Discarded equations create small imbalance}
  \begin{columns}[T,onlytextwidth]
    \begin{column}{0.56\textwidth}
      \centering
      \resizebox{0.93\linewidth}{!}{\StRoundedChords}
    \end{column}
    \begin{column}{0.41\textwidth}
      When a dyadic equation is discarded, fewer than \(L=O(\log n)\) live
      variables can still change.

      \[
        |\Delta(I)|
        \le O\!\left(p_{\max}\log^2(n+1)\right).
      \]

      Boosting gives
      \(p_{\max}=O(1/\log(n+1))\), hence
      \[
        \boxed{|\Delta(I)|=O(\log(n+1)).}
      \]
    \end{column}
  \end{columns}
\end{frame}

% Timing: ~0:55
\begin{frame}{Rounding chooses one chord per client}
  \vspace{0.06cm}
  There exist a backbone tour \(T_R\), auxiliary backbone flow \(y\), and one
  ordered pair \((a_i,b_i)\) for every outside client such that
  {\fontsize{11.0pt}{13.2pt}\selectfont
  \[
    \begin{aligned}
      &\cost(T_R)+\sum_{a,b\in R}d(a,b)y_{ab}\\[-1mm]
      &\qquad+\sum_{i\in U}p_i\bigl(d(a_i,i)+d(i,b_i)\bigr)
      \le O(\log(n+1))\optpost.
    \end{aligned}
  \]}

  \vspace{0.06cm}
  \begin{columns}[T,onlytextwidth]
    \begin{column}{0.30\textwidth}
      \centering
      \stpill[STBlue]{one chord per client}
    \end{column}
    \begin{column}{0.30\textwidth}
      \centering
      \stpill[STTeal]{degree \(O(\log n)\)}
    \end{column}
    \begin{column}{0.30\textwidth}
      \centering
      \stpill[STCoral]{imbalance \(O(\log n)\)}
    \end{column}
  \end{columns}
\end{frame}

% Timing: ~1:00
\begin{frame}{Cut concentration controls every chord set}
  For \(X\subseteq R\), let \(q(X)\) be the number of interval components of
  \(X\) in the cyclic backbone order.

  \vspace{0.16cm}
  With probability at least \(1-(n+1)^{-5}\), simultaneously for every \(X\):
  \[
  \begin{aligned}
    F_A^-(X)&\le2\mu^-(X)+C_0q(X)\log(n+1),\\
    \mu^+(X)&\le2F_A^+(X)+C_0q(X)\log(n+1).
  \end{aligned}
  \]

  \vspace{0.22cm}
  There are at most \((n+1)^{2q}\) subsets with \(q\) interval components, so
  Chernoff bounds survive a union bound over all cuts.
\end{frame}

% Timing: ~1:05
\begin{frame}{Hoffman's condition closes chords into a walk}
  Hoffman's theorem asks, for every \(X\subseteq R\),
  \[
    \ell(\delta^-(X))\le u(\delta^+(X)).
  \]

  \vspace{-0.02cm}
  Use three arc types:
  \begin{center}
  \begin{tabular}{>{\raggedright\arraybackslash}p{0.43\linewidth}
                  >{\centering\arraybackslash}p{0.19\linewidth}
                  >{\raggedright\arraybackslash}p{0.26\linewidth}}
    \toprule
    \textbf{arc type} & \textbf{lower}--\textbf{upper} & \textbf{purpose}\\
    \midrule
    active outside chord & \(1\)--\(4\) & serve every active client\\
    auxiliary arc \(a\to b\) & \(0\)--\(\lceil2y_{ab}\rceil\) & repair mean imbalance\\
    backbone cycle edge & \(1\)--\(O(\log n)\) & connect every cut\\
    \bottomrule
  \end{tabular}
  \end{center}

  \small Interval discrepancy plus cut concentration implies the condition.
\end{frame}

% Timing: ~0:55
\begin{frame}{The realized object has bounded multiplicity}
  \begin{columns}[T,onlytextwidth]
    \begin{column}{0.58\textwidth}
      \centering
      \resizebox{0.91\linewidth}{!}{\StOfflineWalk}
    \end{column}
    \begin{column}{0.39\textwidth}
      With high probability, an integral circulation exists.

      \vspace{0.30cm}
      Expanding every chord
      \(a_i\to b_i\) back to
      \(a_i\to i\to b_i\) yields a balanced, strongly connected multigraph.

      \vspace{0.30cm}
      Its Euler tour is a closed walk serving all active outside clients.
    \end{column}
  \end{columns}
\end{frame}

% Timing: ~1:00
\begin{frame}{The proof ends where causality begins}
  \begin{tikzpicture}[x=1cm,y=1cm]
    \node[draw=STBlue,fill=STBlue!9,rounded corners=0.65pt,
      minimum width=2.15cm,minimum height=0.90cm,
      font=\fontsize{10.5pt}{12pt}\selectfont\bfseries] (backbone) at (1.1,1.55)
      {backbone};
    \node[draw=STTeal,fill=STTeal!9,rounded corners=0.65pt,
      minimum width=2.15cm,minimum height=0.90cm,
      font=\fontsize{10.5pt}{12pt}\selectfont\bfseries] (fractional) at (4.0,1.55)
      {fractional flow};
    \node[draw=STGold,fill=STGold!12,rounded corners=0.65pt,
      minimum width=2.15cm,minimum height=0.90cm,
      font=\fontsize{10.5pt}{12pt}\selectfont\bfseries] (rounded) at (6.9,1.55)
      {rounded chords};
    \node[draw=STCoral,fill=STCoral!9,rounded corners=0.65pt,
      minimum width=2.15cm,minimum height=0.90cm,
      font=\fontsize{10.5pt}{12pt}\selectfont\bfseries] (offline) at (9.8,1.55)
      {offline walk};
    \draw[stArrow] (backbone)--(fractional);
    \draw[stArrow] (fractional)--(rounded);
    \draw[stArrow] (rounded)--(offline);
    \draw[STCoral,line width=1.4pt,densely dashed,
      -{Stealth[length=2.2mm]}] (offline.south)
      to[bend left=28]
      node[below=0.20cm,align=center,text=STCoral,
        font=\fontsize{10pt}{12pt}\selectfont\bfseries]
      {missing: causal,\\cost-sensitive execution}
      (6.7,-0.55);
    \node[draw=STCoral,dashed,rounded corners=0.65pt,
      minimum width=2.40cm,minimum height=0.90cm,
      font=\fontsize{10.5pt}{12pt}\selectfont\bfseries,text=STCoral]
      at (5.25,-0.55) {adaptive policy};
  \end{tikzpicture}

  \vspace{0.08cm}
  \begin{sttightitemize}
    \item the Hoffman circulation depends on the complete active set \(A\);
    \item the next call must be chosen without knowing future active chords;
    \item bounded multiplicity alone does not control the cost of tiny-\(y\)
      auxiliary arcs.
  \end{sttightitemize}
\end{frame}

% ==========================================================================
% DP — 10 minutes
% ==========================================================================
\section{A priori TSP via dynamic programming}
\stsectionframe{STGold}{Part III}{Dynamic programming with\\finite memory}{Survival
suppresses rare long shortcuts;\\a short suffix contains all later terms.}

% Timing: ~1:15
\begin{frame}{A fixed order is a sum of possible shortcuts}
  \centering
  \resizebox{0.83\linewidth}{!}{\StPairwiseShortcut}

  \vspace{0.04cm}
  For \(\pi=(r,v_1,\ldots,v_n,r)\),
  {\fontsize{12.3pt}{14.8pt}\selectfont
  \[
    \cost(\pi)
    =
    \sum_{0\le i<j\le n+1}
      p_i p_j Q(i,j)d(v_i,v_j),
    \qquad
    Q(i,j)=\prod_{i<k<j}(1-p_k).
  \]}
  \small Pair \((i,j)\) is paid exactly when its endpoints are consecutive
  active vertices.
\end{frame}

% Timing: ~1:30
\begin{frame}{Extra activation damps rare long gaps}
  \centering
  \resizebox{0.83\linewidth}{!}{\StActivationInflation}

  \vspace{0.04cm}
  Retain pairs with \(Q(i,j)>\eta\). For
  \(\widetilde p_v=1-(1-p_v)^{1+\alpha}\) and
  \(\kappa=(1+\alpha)^2\),
  \[
    \cost(\pi)
    \le
    \frac{\kappa}{1-\kappa\eta^\alpha}\,
    \cost^{\mathrm{loc}}_\eta(\pi)
    \qquad(\kappa\eta^\alpha<1).
  \]
\end{frame}

% Timing: ~1:10
\begin{frame}{Blocks make the interface finite}
  \centering
  \resizebox{0.92\linewidth}{!}{\StBlockStrip}

  \vspace{0.12cm}
  Fix \(0<\delta<1/4\):
  \begin{itemize}
    \item a \textbf{heavy} client has \(p_v>\delta/2\) and forms a singleton;
    \item a \textbf{light block} \(B\) has
      \(\mu(B)=\sum_{v\in B}p_v\le\delta\);
    \item after the first light block in a run,
      \(\mu(B)\ge\delta/2\).
  \end{itemize}
  Scan each light run from right to left to obtain such a consecutive
  partition.
\end{frame}

% Timing: ~1:20
\begin{frame}{The block objective has controlled loss}
  For a block \(B\), precompute its best internal order:
  \[
    I(B)=\min_\sigma\sum_{x<_\sigma y}p_xp_y d(x,y).
  \]
  For an ordered partition \(\mathcal B=(B_1,\ldots,B_m)\),
  {\fontsize{11.6pt}{14pt}\selectfont
  \[
    F_\eta(\mathcal B)
    =
    \sum_t I(B_t)
    +
    \sum_{\substack{i<j\\ \prod_{i<t<j}q(B_t)>\eta}}
      \left(\prod_{i<t<j}q(B_t)\right)C(B_i,B_j).
  \]}

  \vspace{0.02cm}
  \[
    \cost^{\mathrm{loc}}_\eta(\pi(\mathcal B))
    \le F_\eta(\mathcal B),
    \qquad
    F_\eta(\mathcal B_\pi)
    \le(1-\delta)^{-2}\cost(\pi).
  \]
  \vspace{-0.08cm}
  \small Only endpoint-internal survival factors are dropped; each is at
  least \(1-\delta\).
\end{frame}

% Timing: ~1:20
\begin{frame}{Only a short suffix can affect the future}
  \centering
  \resizebox{0.84\linewidth}{!}{\StDPSuffix}

  \vspace{0.04cm}
  A state stores
  \[
    (U,\Omega)
    =
    \bigl(\text{placed clients},\ \text{relevant suffix of blocks}\bigr).
  \]
  \small Once intervening survival falls to at most \(\eta\), an old block
  cannot interact with any future block.
\end{frame}

% Timing: ~1:25
\begin{frame}{Appending a block exposes its local terms}
  If the relevant suffix is
  \(\Omega=(B_a,\ldots,B_t)\), appending \(B\) adds
  \[
    I(B)
    +
    \sum_{s=a}^{t}
      \left(\prod_{s<u\le t}q(B_u)\right)C(B_s,B).
  \]

  \vspace{0.08cm}
  Then trim the suffix by the same product rule.

  \vspace{0.10cm}
  \begin{columns}[T,onlytextwidth]
    \begin{column}{0.48\textwidth}
      {\color{STBlue}\rule{3pt}{2.3em}}\hspace{0.35em}%
      \parbox[c]{0.88\linewidth}{\textbf{Exactness.} These are exactly the
      terms created by appending \(B\).}
    \end{column}
    \begin{column}{0.48\textwidth}
      {\color{STTeal}\rule{3pt}{2.3em}}\hspace{0.35em}%
      \parbox[c]{0.88\linewidth}{\textbf{Memory.} The suffix contains
      \(K=O((\log(1/\eta)+1)/\delta)\) blocks.}
    \end{column}
  \end{columns}

  \vspace{0.10cm}
  The resulting shortest-path DP runs in
  \(c_{\delta,\eta}^{\,n}\poly(n)\).
\end{frame}

% Timing: ~1:40
\begin{frame}{The losses multiply to at most \(1+\varepsilon\)}
  Choose
  \[
    \delta=\frac{\varepsilon}{16},
    \qquad
    \alpha=\frac{\varepsilon}{32},
    \qquad
    \eta=\left(\frac{\varepsilon}{64}\right)^{32/\varepsilon}.
  \]
  Then \(\eta^\alpha=\varepsilon/64\), and the DP output
  \(\widehat\pi\) satisfies
  {\fontsize{13.4pt}{16pt}\selectfont
  \[
  \begin{aligned}
    \cost(\widehat\pi)
      &\le
      \frac{(1+\alpha)^2}
           {1-(1+\alpha)^2\eta^\alpha}
      F_\eta(\widehat{\mathcal B})\\
      &\le
      \frac{(1+\alpha)^2}
           {1-(1+\alpha)^2\eta^\alpha}
      (1-\delta)^{-2}\optprior\\
      &\le(1+\varepsilon)\optprior.
  \end{aligned}
  \]}

  \centering
  \stpill[STGold]{deterministic runtime \(c_\varepsilon^n\poly(n)\)}
\end{frame}

% ==========================================================================
% Close
% ==========================================================================
\section*{Takeaway}
\begin{frame}[plain]
  \begin{tikzpicture}[remember picture,overlay]
    \fill[STPaper] (current page.south west) rectangle (current page.north east);
    \draw[STLine,line width=0.45pt]
      ([xshift=0.95cm,yshift=-0.82cm]current page.north west)
      --
      ([xshift=-0.88cm,yshift=-0.82cm]current page.north east);
    \draw[STTeal,line width=1.35pt]
      ([xshift=0.95cm,yshift=-0.82cm]current page.north west)
      --
      ([xshift=3.15cm,yshift=-0.82cm]current page.north west);
    \node[
      anchor=west,
      text=STTeal,
      font=\fontsize{9pt}{11pt}\selectfont\bfseries
    ] at ([xshift=0.98cm,yshift=2.42cm]current page.west)
      {TAKEAWAY};
    \node[
      anchor=west,
      text=STInk,
      text width=11.4cm,
      align=left,
      font=\fontsize{21.5pt}{25.5pt}\selectfont\bfseries
    ] at ([xshift=0.98cm,yshift=1.42cm]current page.west)
      {Information is a resource; causality is the constraint.};

    \draw[STBlue,line width=1.1pt]
      ([xshift=0.98cm,yshift=0.58cm]current page.west)
      --
      ([xshift=3.88cm,yshift=0.58cm]current page.west);
    \draw[STTeal,line width=1.1pt]
      ([xshift=4.72cm,yshift=0.58cm]current page.west)
      --
      ([xshift=7.62cm,yshift=0.58cm]current page.west);
    \draw[STGold,line width=1.1pt]
      ([xshift=8.46cm,yshift=0.58cm]current page.west)
      --
      ([xshift=11.36cm,yshift=0.58cm]current page.west);

    \node[anchor=north west,text=STInkSoft,text width=3.25cm,align=left,
      font=\fontsize{11.1pt}{14.2pt}\selectfont]
      at ([xshift=0.98cm,yshift=0.34cm]current page.west)
      {\hyphenpenalty=10000
       \textbf{\color{STBlue}Gap constructions}\par
       \vspace{0.12cm}
       Three distinct bottlenecks:\par
       \(\tfrac54,\ \tfrac{11}{10},\ \tfrac43\).};
    \node[anchor=north west,text=STInkSoft,text width=3.25cm,align=left,
      font=\fontsize{11.1pt}{14.2pt}\selectfont]
      at ([xshift=4.72cm,yshift=0.34cm]current page.west)
      {\hyphenpenalty=10000
       \textbf{\color{STTeal}Circulation}\par
       \vspace{0.12cm}
       Offline structure exists; the causal step remains open.};
    \node[anchor=north west,text=STInkSoft,text width=3.78cm,align=left,
      font=\fontsize{11.1pt}{14.2pt}\selectfont]
      at ([xshift=8.46cm,yshift=0.34cm]current page.west)
      {\hyphenpenalty=10000
       \textbf{\color{STGold}A priori DP}\par
       \vspace{0.12cm}
       Survival locality gives a \((1+\varepsilon)\)-approximation in
       \(c_\varepsilon^n\poly(n)\).};

    \draw[STLine,line width=0.45pt]
      ([xshift=0.98cm,yshift=-2.25cm]current page.west)
      --
      ([xshift=-0.90cm,yshift=-2.25cm]current page.east);
    \node[
      anchor=west,
      text=STInk,
      text width=10.9cm,
      font=\fontsize{11.8pt}{14.5pt}\selectfont\bfseries
    ] at ([xshift=0.98cm,yshift=-2.86cm]current page.west)
      {Open question: can the offline circulation be made causal?};
    \node[
      anchor=east,
      text=STMuted,
      font=\fontsize{8.5pt}{10pt}\selectfont\bfseries
    ] at ([xshift=-0.90cm,yshift=-2.86cm]current page.east)
      {QUESTIONS};
  \end{tikzpicture}
\end{frame}

\end{document}
